\documentclass[12pt]{report}
\usepackage[utf8]{inputenc}
\usepackage[T1]{fontenc}
\usepackage{fontspec}
\setmainfont{Times New Roman}
\newfontfamily\hackfont{Hack}
\usepackage{titlesec}
\usepackage{fancyhdr}
\usepackage{listings}
\usepackage{xcolor}
\usepackage{graphicx}
\usepackage{longtable}
\usepackage{caption}
\usepackage{hyperref}
\usepackage{setspace}
\usepackage{geometry}
\geometry{a4paper, margin=1in}
\hypersetup{colorlinks=true, linkcolor=blue, urlcolor=blue}
\renewcommand{\contentsname}{Table of Contents}

\pagestyle{fancy}
\fancyhf{} 							% Header Middle
\rhead{VLSI Design Lab}
\lhead{Lab-04: FlipFlop Circuit}
\rfoot{\thepage}

\begin{document}

\begin{titlepage}					% Cover Start
\begin{center}
    
\includegraphics[width=0.35\textwidth]{logo.png} \\
    \vspace{0.5cm}
    \textbf{\Huge \textcolor{cyan}{NOTRE DAME} \textcolor{cyan}{UNIVERSITY}} \\
    \vspace{0.5cm}
    \textbf{\Huge \textcolor{orange}{BANGLADESH}} \\
    
    \vspace{0.9cm}
    \textbf{\Huge \underline{VLSI Design Lab Report}} \\
    
    \vspace{1.5em}
    \begin{flushleft}
    	\textbf{\Large Course Code: CSE-4218} \\
		\vspace{0.3cm}        
        \textbf{\Large Course Title: VLSI Design Lab} \\
        \vspace{0.3cm}
        \textbf{\Large Experiment Name: } \\        
    \end{flushleft}
    
    \vspace{1em}
    \begin{flushleft}
        \textbf{\Huge \textcolor{blue}{Submitted by:}} \\
        \vspace{0.3cm}
        \textbf{\Large Name: } \\
        \vspace{0.3cm}
        \textbf{\Large ID: } \\
		\vspace{0.3cm}        
        \textbf{\Large Batch: CSE-20} \\
        \vspace{0.5cm}
        \textbf{\Large Submission Date: }{\Large \textbf{\today}\par}
    \end{flushleft}
    \vfill
    \begin{flushleft}
        \textbf{\Huge \textcolor{blue}{Submitted to:}} \\
        \vspace{0.3cm}
        \textbf{\Large Nafisa Tabassum,} \\
        \vspace{0.3cm}
        \textbf{\Large Lecturer, Dept. of CSE} \\
        \vspace{0.3cm}
        \textbf{\Large Notre Dame University Bangladesh}
    \end{flushleft}
\end{center}
\end{titlepage}						% Cover End

\tableofcontents
\thispagestyle{empty}
\clearpage
\pagenumbering{arabic}

\section*{\center Abstract (All Lab)}


\newpage
\chapter{Lab-01 CMOS Logic Gate}
\section{Introduction}


\section{Objective}


\section{Function}



\subsection*{Truth Table}
\begin{center}
\begin{tabular}{|c|c|c|}
\hline
Clock & D & Q(next) \\ \hline
$\uparrow$ & 0 & 0 \\
$\uparrow$ & 1 & 1 \\
No Edge & X & Q(previous) \\ \hline
\end{tabular}
\end{center}

\subsection*{Schematic Image}
\begin{figure}[!h]
\centering
%\includegraphics[width=1\textwidth]{FlipFlop.png}
\caption{CMOS Schematic of D Flip-Flop in DSCH2}
\end{figure}


\subsection*{Timing Diagram}
\begin{figure}[!h]
\centering
%\includegraphics[width=1\textwidth]{Timing Diagram.png}
\caption{Timing Diagram of D Flip-Flop}
\end{figure}

The timing diagram shows that the output Q follows the input D only at the rising edge of the clock signal. Between clock edges, the output remains constant regardless of changes in the input. \\
From the simulation results, it was observed that the D Flip-Flop changes its output state only at the rising edge of the clock signal. When the clock transitions from LOW to HIGH, the output Q takes the value of the input D. At all other times, the output remains unchanged, preserving the previously stored value. This behavior confirms the correct operation of the edge-triggered D Flip-Flop implemented using CMOS logic.

% =====================================================
\section{Conclusion}
The CMOS-based D Flip-Flop was successfully designed and simulated using the DSCH2 tool. The truth table and timing diagram verified that the circuit operates correctly as a clocked storage element. This experiment reinforces the understanding of sequential logic design and the role of flip-flops in digital systems.



\chapter{Lab-02 Sequential Circuit}



\chapter{Lab-03}


\chapter{Lab-04 }

\chapter{Lab-05 }

\chapter{Lab-06 }




\end{document}